\section{Strips}\label{sec:strip}

On two rows every rectangle is axis-aligned. Indeed, no side of a tilted rectangle
is horizontal, since the perpendicular side would then be vertical. Let $p$ be a
lowest corner, with sides $u$ and $v$ at $p$. Then $p+u$ and $p+v$ lie strictly
above $p$, and $p+u+v$ lies strictly above both. So a tilted rectangle needs three
rows.
An axis-aligned rectangle on a strip uses both rows and two columns $a<b$, and we
call it $R_{a,b}$.

\begin{theorem}[Strip Theorem]\label{thm:strip}
Let $m\ge3$.
\begin{enumerate}[(a)]
\item A blackout of the $2\times m$ strip is valid if and only if at most one of its
columns is full.
\item The maximum blackout is $f(2,m)=m+1$, and the maximum blackouts are exactly the
blackouts with one full column and exactly one black point in every other column.
\item There are $a_2(m)=m\cdot2^{m-1}$ maximum blackouts.
\end{enumerate}
For $m=2$ there is only one rectangle, so every blackout is valid, $f(2,2)=4$ and
$a_2(2)=1$.
\end{theorem}

\begin{proof}
\emph{Step 1. The constraints.} Two distinct rectangles $R_{a,b}$ and $R_{c,d}$
cannot share both columns. If they share one column, that column cancels in the
difference set. What remains is the four points of the other two columns $p$ and
$q$, and we call it the \emph{pair constraint} on $\{p,q\}$. If they share no
column, the difference set is all eight points of four columns, and it contains a
pair constraint. Since $m\ge3$ there is a third column $z$. The rectangles on the
column pairs $\{z,p\}$ and $\{z,q\}$ have difference set exactly the pair
constraint on $\{p,q\}$, so every pair constraint occurs. By Corollary~\ref{cor:hitting} a
blackout $S$ is valid if and only if the kept set $K$ meets every pair of columns.
That is, $S$ is valid if and only if at most one column has no kept point. This is
(a).

\emph{Step 2. The bound.} Let $S$ be valid. By (a), at least $m-1$ columns contain a
kept point, so $|K|\ge m-1$ and $|S|=2m-|K|\le m+1$. If $|S|=m+1$, then $|K|=m-1$,
so exactly one column is full and every other column has exactly one kept point,
hence exactly one black point. Conversely, such a blackout has one full column, so
it is valid by (a), and it has $m+1$ points. This is (b).

\emph{Step 3. The count.} There are $m$ choices for the full column and two choices
in each of the other $m-1$ columns. This is (c).
\end{proof}

Figure~\ref{fig:stripmax} shows one of the $6\cdot2^5=192$ maximum blackouts of the
$2\times6$ strip. The count $80$ for the $2\times5$ strip was written in the ledger
before the program first ran on that grid, and the program found $80$.
Row $2$ of \seqnum{A399580} is $12,32,80,192,\ldots$ from $m=3$ on, which is
\seqnum{A001787}$(m)=m2^{m-1}$, and only $a_2(2)=1$ differs from
$\mathrm{A001787}(2)=4$.

\begin{figure}[t]
\centering
\begin{tikzpicture}[x=0.700cm,y=0.700cm,baseline=(current bounding box.north)]
  \fill[keptcol] (3.6,-0.6) rectangle (4.4,-2.4);
  \draw[gridcol] (1,-1) grid (6,-2);
  \fill[black] (1,-1) circle (0.17);
  \filldraw[fill=white,draw=keptline,line width=0.6pt] (2,-1) circle (0.17);
  \filldraw[fill=white,draw=keptline,line width=0.6pt] (3,-1) circle (0.17);
  \fill[black] (4,-1) circle (0.17);
  \fill[black] (5,-1) circle (0.17);
  \filldraw[fill=white,draw=keptline,line width=0.6pt] (6,-1) circle (0.17);
  \filldraw[fill=white,draw=keptline,line width=0.6pt] (1,-2) circle (0.17);
  \fill[black] (2,-2) circle (0.17);
  \fill[black] (3,-2) circle (0.17);
  \fill[black] (4,-2) circle (0.17);
  \filldraw[fill=white,draw=keptline,line width=0.6pt] (5,-2) circle (0.17);
  \fill[black] (6,-2) circle (0.17);
  \node[lab] at (1,-0.35) {1};
  \node[lab] at (2,-0.35) {2};
  \node[lab] at (3,-0.35) {3};
  \node[lab] at (4,-0.35) {4};
  \node[lab] at (5,-0.35) {5};
  \node[lab] at (6,-0.35) {6};
  \node[lab] at (0.35,-1) {1};
  \node[lab] at (0.35,-2) {2};
\end{tikzpicture}

\caption{A maximum blackout of the $2\times6$ strip. Column $4$ (shaded) is full and
every other column has one black point, which gives $m+1=7$ points.}
\label{fig:stripmax}
\end{figure}

Part (a) of the theorem also counts the valid blackouts of a strip by size.

\begin{proposition}[size polynomial]\label{prop:stripsize}
For $m\ge3$ let $F_m(u)=\sum_Su^{|S|}$, summed over the valid blackouts $S$ of the
$2\times m$ strip. Then
\[
F_m(u)=(1+2u)^m+m\,u^2(1+2u)^{m-1}.
\]
\end{proposition}

\begin{proof}
By Theorem~\ref{thm:strip}(a) a blackout is valid if and only if at most one column
is full. A column that is not full is empty or has exactly one of its two points
black, so it contributes the factor $1+2u$. The first term counts the blackouts with
no full column. In the second term the full column is chosen in $m$ ways and
contributes $u^2$, and the other $m-1$ columns are not full.
\end{proof}

The formula fails at $m=2$, where the all-black board is valid and $F_2(u)=(1+u)^4$.

For $m\ge3$ the strip is not a square. So its symmetries form the group of order
four generated by the row reflection $j\mapsto3-j$ and the column reversal
$i\mapsto m+1-i$.

\begin{proposition}[orbits]\label{prop:striporbits}
For $m\ge3$ the number of maximum blackouts of the $2\times m$ strip up to the
symmetries of the strip is
\[
o_m=m\,2^{m-3}+[m\text{ odd}]\,2^{(m-3)/2}.
\]
\end{proposition}

\begin{proof}
By Burnside's lemma, $4o_m$ is the sum over the four symmetries of the number of
maximum blackouts each one fixes. The identity fixes all $m2^{m-1}$. The row
reflection moves the single black point of every column that is not full, so it
fixes none. The column reversal fixes a maximum blackout
exactly when two things hold. The full column $e$ satisfies $e=m+1-e$, which needs
$m$ odd and $e$ central. And the columns $i$ and $m+1-i$ carry their black points in
the same row, which leaves $2^{(m-1)/2}$ choices. The half turn, the product of the
two, needs the same full column and opposite rows in the columns $i$ and $m+1-i$,
again $2^{(m-1)/2}$ choices. Adding up,
$4o_m=m2^{m-1}+2\cdot[m\text{ odd}]2^{(m-1)/2}$.
\end{proof}

\begin{remark}[nonattacking kings]\label{rem:kings}
The orbit counts $o_m$ for $m=3,4,\ldots$ are $4,8,22,48,116,256,\ldots$. These are
the terms \seqnum{A322284}$(m-1)$, the numbers of ways to place $n=m-1$
nonattacking kings on a $2\times2n$ chessboard up to the symmetries of the board.
This is not a coincidence. Cut the board into $n$ blocks of two adjacent columns.
Two squares of one block are always adjacent, so each block holds exactly one of the
$n$ kings. A king in the right column of a block
attacks every square of the left column of the next block. So once a block uses its
right column, every later block does too. Hence a placement is a cut
$k\in\{0,\dots,n\}$, the number of blocks using their left column, together with a
row for each block. Send it to the strip of width $m=n+1$ whose full column is
column $k+1$. The other $n$ columns carry, in order, the rows of the $n$ blocks.
This is a bijection onto the maximum blackouts. It sends the row reflection
of the board to the row reflection of the strip. The reversal of the board reverses
the order of the blocks and swaps left and right inside each block, so it replaces
$k$ by $n-k$. Hence it goes to the column reversal of the strip. So it matches orbits
with orbits, and $o_m=\mathrm{A322284}(m-1)$ for every $m\ge3$. The formula given in
that entry agrees with Proposition~\ref{prop:striporbits}.
\end{remark}
